% Bilan énergétique de l'expérience de Joule pour 20 chutes (corrigé du défi D19)
\begin{tikzpicture}[x=1cm,y=1cm]
\node[draw=otNavy,line width=0.8pt,fill=pcEau,rounded corners=2mm,minimum width=3.2cm,minimum height=1.3cm,align=center,pcLegendeG] (M) at (0,0) {masses qui tombent\\ $\Delta E_\mathrm{pp}=\qty{-8,16}{kJ}$};
\node[draw=otNavy,line width=0.8pt,fill=pcEau,rounded corners=2mm,minimum width=3.6cm,minimum height=1.3cm,align=center,pcLegendeG] (E) at (6.6,0) {eau, récipient, palettes\\ $\Delta U=\qty{7,98}{kJ}$ (mesurée)};
\draw[otNavy,line width=1.6pt,-{Stealth[length=3mm]}] (M.east) -- (E.west) node[pos=0.5,above,pcLegende,text=otNavy] {travail $W=\qty{8,16}{kJ}$};
\draw[otOrange!90!black,line width=1.1pt,-{Stealth[length=2.4mm]},decorate,
  decoration={snake,amplitude=0.8mm,segment length=3.5mm,post length=2.5mm}] (E.south) -- ++(0,-1.3)
  node[anchor=north,pcLegende,text=otOrange!80!black,align=center] {pertes : environ \qty{0,19}{kJ}\\ (air, poulies, frottements)};
\end{tikzpicture}
\par\smallskip\legende{Presque tout le travail des masses se retrouve en énergie interne de l'eau ; les pertes, d'environ 2~\%, faussent légèrement la mesure.}
